\section{Technical implementation}
\subsection{Topology and treatments}
The intended physical path is client MacBook, Linux NucBox router, then server
MacBook, with dedicated wired links and separate Wi-Fi management. The client
uses 10.203.0.2/30 and the server uses 10.203.0.6/30. Each has an explicit host
route to its peer through the router. Our preflight refuses a selected Mac
default-route interface, incorrect address or mask, or mismatched source revision.
Each workload also requires an active Ethernet link after router setup.
The Mac helpers inspect existing network settings
without changing them. The Linux helper owns a private router namespace and
temporarily attaches only the two explicitly selected, unused Ethernet ports.

The shaped return path is 100 Mb/s. FIFO uses a 250,000-byte queue below HTB;
FQ-CoDel uses the same shaper with its own queue. The CAKE treatments shape
directly at 100 Mb/s. The bounded variant classifies only the designated
foreground reply flow and applies a 5 Mb/s policing budget. Incoming DS fields
are cleared before our own classification. These are documented candidate
settings that still need team acceptance on the final hardware.

In the mixed topology, NetEm adds 20 ms on router egress toward the server.
The Linux-only topology instead adds 10 ms on each endpoint's ingress.
Both target about 20 ms of added base RTT, but the queue placement differs.
We record the mixed topology separately and do not pool its results with the
namespace or all-Linux physical runs. The Mac endpoints retain native TCP
behavior. Linux TCP\_INFO and AF\_PACKET measurements are unavailable there.

\subsection{Workloads and lifecycle}
The candidate foreground workload sends 512-byte messages at 500 requests per
second using the client's monotonic clock. UDP is the main planned transport;
TCP is a separately labeled sensitivity experiment with fixed framing and
TCP\_NODELAY. Reverse TCP iperf3 traffic supplies background demand at 50\%,
90\% or 110\% of the shaped rate, split across one or eight flows. Offered
demand above capacity does not imply delivered throughput above capacity.

The pilot has two blocks, four policies, unloaded baselines and a 110 Mb/s
single-flow condition: 16 trials. The larger plan has ten 60-second trials
for each loaded condition and ten unloaded trials per policy: 280 trials.
Policy order is randomized within recorded blocks. Physical collection is
serial because simultaneous trials would share the bottleneck and invalidate
the comparison. Local simulation uses up to eight worker processes; its
default six blocks produce 48 trials and 240,000 scheduled slots, three times
the pilot volume.

Each trial records its source, configuration, lifecycle, observations and
artifact hashes in a fresh directory. Failed attempts remain available.
Commands have deadlines and bounded output. Cancellation stops owned child
processes and preserves cleanup errors. Endpoint leases prevent overlapping
trials. Linux cleanup verifies resource identity before returning interfaces.
The pipeline proceeds automatically from successful collection to checked
export, analysis, plots and a text interpretation.

\subsection{Analysis and verification}
We calculate per-run latency percentiles and match loaded runs to their
unloaded baseline in the same block. Added p99 is loaded p99 minus unloaded
p99. Policy comparisons use paired block resampling for 95\% bootstrap
intervals; individual requests are not independent replications. The pipeline
retains generator omissions, timeouts and late replies. A low on-time p99
without those outcomes can be misleading.

Our exploratory analysis standardizes six run metrics and fits K-means, Ward
agglomerative clustering and DBSCAN. PCA and t-SNE display those metrics;
t-SNE itself is not a clustering algorithm. Noise points remain in all
comparisons. Ridge and Lasso use design features with entire blocks held out;
scaling and penalty selection use only training folds. At least six blocks
are required. These models describe observed designs and do not establish
causality or transfer to new hardware.

Jonathan's original comparison workflow and Chris's C++ networking prototype
remain in the repository with their authorship intact. The C++ probe is a
separate workload check, not the source of the main foreground metric. Our
unit and integration checks cover data validation, real sockets, subprocess
failures and cleanup. CI enforces 95\% statement and branch coverage separately
for core and analysis, with a separate Linux job for privileged testbed checks.
Hardware acceptance remains a separate activity.

For $N$ requests, per-run sorting costs $O(N\log N)$ time and $O(N)$ auxiliary
memory. Analysis bounds each export batch at 200,000 request rows and keeps
at most 400 run summaries for the offline models. Ward and DBSCAN can require
$O(R^2)$ work or storage for $R$ runs; the fixed run cap bounds that cost.
Foreground accounting bounds active requests at 1,024. Ruff enforces McCabe
complexity at most 10 per function. These limits support review and predictable
failure behavior; coverage and complexity scores do not certify correctness.
